\addcontentsline{toc}{subsubsection}{Solution 17. IMO 2022 Shortlist -- Algebra A3}

\begin{dlsummary}[title={\bfseries Solution 17. IMO 2022 Shortlist -- Algebra A3},breakable]

\begin{proof}
Call a pair $(x,y)$ \emph{good} if
\[
xf(y)+yf(x)\leq2.
\]

\medskip
\noindent
\textbf{Step 1. Every good pair has $x=y$.}

Suppose $(x,y)$ is good with $x\ne y$.
By symmetry and uniqueness, neither
$(x,x)$ nor $(y,y)$ can be good. Thus
\[
xf(x)>1,\qquad yf(y)>1.
\]

By AM--GM,
\[
\begin{aligned}
xf(y)+yf(x)
&\geq2\sqrt{xf(y)\cdot yf(x)}\\
&=2\sqrt{xf(x)\cdot yf(y)}\\
&>2,
\end{aligned}
\]
a contradiction.

Therefore every good pair satisfies $x=y$.

\medskip
\noindent
\textbf{Step 2. Establish an upper bound.}

For every $x>0$, a good pair exists.
By Step 1, it must be $(x,x)$. Hence
\[
2xf(x)\leq2,
\]
so
\[
\boxed{f(x)\leq\frac1x.} \tag{1}
\]

\medskip
\noindent
\textbf{Step 3. Force equality.}

Fix $x>0$ and let
\[
t=\frac1{f(x)}>0.
\]

By (1),
\[
f(t)\leq\frac1t=f(x).
\]
Therefore,
\[
\begin{aligned}
xf(t)+tf(x)
&\leq xf(x)+1\\
&\leq2.
\end{aligned}
\]

Thus $(x,t)$ is a good pair. By Step 1,
$t=x$, which gives
\[
\boxed{f(x)=\frac1x\qquad(x>0).}
\]

\medskip
\noindent
\textbf{Verification.}

For $f(x)=1/x$, the given inequality becomes
\[
\frac{x}{y}+\frac{y}{x}\leq2.
\]

However, AM--GM gives
\[
\frac{x}{y}+\frac{y}{x}\geq2,
\]
with equality if and only if $x=y$.

Hence, for every $x>0$, there is exactly
one admissible $y>0$, namely $y=x$.

Therefore $f(x)=1/x$ is the unique solution.
\end{proof}

\medskip
\noindent
\textbf{Alternative 1.1. A direct contradiction.}

We use the bound $f(x)\leq1/x$ established above. Suppose that
$f(x)<1/x$ for some $x>0$. Then the interval $(x,1/f(x))$ is nonempty.
For every $a$ in this interval,
\[
af(x)<1,\qquad xf(a)\leq\frac{x}{a}<1,
\]
so
\[
af(x)+xf(a)<2.
\]
Thus every such $a$ is a good partner of $x$. Since the interval contains
more than one number, this contradicts uniqueness. Hence $f(x)=1/x$.

\medskip
\noindent
\textbf{Alternative 1.2. Strict decrease.}

We again use the fact that every good pair is diagonal and
$f(x)\leq1/x$. If $y>x$, then $(x,y)$ is not good, whereas $(y,y)$ is
good. Therefore
\[
xf(y)+yf(x)>2\geq2yf(y)>xf(y)+yf(y),
\]
which gives $f(x)>f(y)$. Thus $f$ is strictly decreasing.

For any $y>x$, strict decrease and the fact that $(x,y)$ is not good give
\[
2<xf(y)+yf(x)<(x+y)f(x).
\]
Letting $y\to x^+$ yields $f(x)\geq1/x$. Together with the upper bound,
this proves $f(x)=1/x$.

\medskip
\noindent
\textbf{Alternative 2. A partner map.}

For each $x>0$, let $g(x)$ be the unique good partner of $x$. Symmetry of
the condition and uniqueness imply $g(g(x))=x$.

We first show that $f$ is strictly decreasing. If $x<y$, then
$(g(y),y)$ is good, but $(g(y),x)$ is not, since $y$ is the unique good
partner of $g(y)$. Hence
\[
xf(g(y))+g(y)f(x)>2\geq yf(g(y))+g(y)f(y).
\]
Rearranging shows $f(x)>f(y)$.

Now suppose that $(x,y)$ is a good pair with $x\ne y$. Uniqueness
implies that neither $(x,x)$ nor $(y,y)$ is good, so
$xf(x)+yf(y)>2$. But $(x,y)$ being good gives
\[
(x-y)(f(x)-f(y))
=xf(x)+yf(y)-xf(y)-yf(x)>0,
\]
contradicting strict decrease. Thus every good pair is diagonal, giving
$f(x)\leq1/x$. Alternative 1.1 now forces equality.

\medskip
\noindent
\textbf{Alternative 3. Equality at the unique partner.}

Let $g(x)$ again denote the unique good partner of $x$. Symmetry gives
$g(g(x))=x$. We claim that every good pair satisfies equality. Suppose
instead that $(x,y)$ is good and
\[
\delta=2-xf(y)-yf(x)>0.
\]
For $0<\varepsilon<\delta/(2f(x))$, uniqueness makes $(x,y+\varepsilon)$
not good, and therefore
\[
f(y+\varepsilon)
>\frac{2-(y+\varepsilon)f(x)}{x}
>f(y)+\frac{\delta}{2x}.
\]
Put $z=y+\varepsilon$. The good pair $(z,g(z))$ and the fact that
$(g(z),y)$ is not good give
\[
f(g(z))\leq\frac{2-g(z)f(z)}{z},
\qquad
f(g(z))>\frac{2-g(z)f(y)}{y}.
\]
Combining these inequalities yields
\[
2\varepsilon<g(z)\bigl[zf(y)-yf(z)\bigr]
<g(z)\left[\varepsilon f(y)-\frac{y\delta}{2x}\right].
\]
For sufficiently small positive $\varepsilon$, the final bracket is
negative, a contradiction. Thus every good pair has value exactly $2$.

Every pair other than $(x,g(x))$ is not good, so
\[
xf(y)+yf(x)\geq2\qquad(x,y>0).
\]
Taking $y=x$ gives $f(x)\geq1/x$. Applying this bound to the good pair
$(x,g(x))$ gives
\[
2\geq xf(g(x))+g(x)f(x)
\geq\frac{x}{g(x)}+\frac{g(x)}{x}\geq2.
\]
Equality throughout implies $f(x)=1/x$.

\medskip
\noindent
\textbf{Alternative 4. Strict convexity and continuity.}

Let $g(w)$ be the unique good partner of $w$, and define
\[
q_s(t)=f(t)+st.
\]
For $s=f(g(w))/g(w)$, the good-pair condition gives
$q_s(w)\leq2/g(w)$. If $t\ne w$, then $(g(w),t)$ is not good, so
$q_s(t)>2/g(w)$. Thus $q_s$ has a unique global minimum at $w$.

We claim that $f$ is strictly convex. Otherwise, for some $u<v$ and
$\lambda\in(0,1)$, setting $w=\lambda u+(1-\lambda)v$ would give
\[
f(w)\geq\lambda f(u)+(1-\lambda)f(v).
\]
Adding the linear term $st$ yields
$q_s(w)\geq\lambda q_s(u)+(1-\lambda)q_s(v)$, contradicting the
strict minimality of $q_s$ at $w$. Hence $f$ is strictly convex and, as
a finite convex function on an open interval, is continuous.

As in the first solution, every good pair is diagonal, so $f(x)\leq1/x$.
If this inequality were strict at some $x$, the continuous function
$h(y)=xf(y)+yf(x)$ would satisfy $h(x)=2xf(x)<2$. It would then be less
than $2$ for all $y$ in some interval around $x$, contradicting the
uniqueness of the good partner. Therefore $f(x)=1/x$ for all $x>0$.

\end{dlsummary}